\chapter{Architectural Redirections}

This appendix records the explicit design decisions that changed the project's
direction from its earlier experiments. It is drawn from the architecture
document \texttt{docs/charlotte-sitas-xous-architecture.md} and reflects the
transition from an earlier ``universal completion-capability'' model to the
three-primitive architecture described in this manual.

\section{What to retain from the current implementation}

The following work aligns with the architecture and should continue:

\begin{itemize}
    \item AArch64 EL0/SVC path.
    \item Caller identity derived from the current trap/thread context.
    \item Per-address-space capability tables.
    \item Shared CQ ring.
    \item Non-lossy kernel CQ backlog.
    \item \texttt{CQ\_WAIT} blocking through scheduler/observer machinery.
    \item Bounded queues and explicit \texttt{WouldBlock}.
    \item No-std sitas split and CharlotteOS backend.
    \item Per-shard executor model.
    \item Typed in-kernel \texttt{ShardMailbox<M>}.
    \item Closure-based \texttt{ShardLocal<T>} with restricted use.
    \item Segment-aware ELF loading.
    \item QEMU boot CI across AArch64 and x86-64.
\end{itemize}

\section{What to reframe}

These components are correct but their \emph{role} in the architecture has
shifted:

\begin{itemize}
    \item \textbf{Completion capability work} is the \textbf{operation completion
        subsystem}, not the universal service IPC subsystem. Use completion
        queues for kernel and device operations; use endpoints for cross-domain
        service calls.

    \item \textbf{Observer/waker code} is internal scheduling machinery, not the
        public semantic identity of all async objects. A Waker is a hint to
        repoll a future; it is not a completion record.

    \item \textbf{IPI queue closures} are a temporary kernel work mechanism, not
        the long-term cross-domain messaging model. Use typed mailboxes or
        closed kernel enums for subsystem-specific work.

    \item \textbf{The earlier mailbox syscall ABI} is a smoke-test interface
        only. The stable ABI is endpoint IPC.
\end{itemize}

\section{What to replace}

\subsection{Stop expanding raw mailbox syscalls}

A mailbox tied directly to LP identity is not a complete userspace service
abstraction: it exposes placement where authority should be exposed, lacks
protocol identity, conflates shard transport with domain IPC, and lacks
move/lend semantics. \textbf{Direction:} implement endpoint and connection
capabilities independent of LP identity. Clients address the service
connection, not a target LP.

\subsection{Do not allocate one completion capability per service request}

For userspace service calls, the natural object is a reply token. For
high-rate kernel operations, per-operation capabilities create unnecessary
table pressure. \textbf{Direction:} use reply tokens for endpoint calls;
operation IDs plus CQ entries for ordinary kernel operations; first-class
operation capabilities only where independent authority is required.

\subsection{Separate readiness, notification, and completion}

Their contracts differ: readiness means progress may be possible; notification
means inspect state; completion means an operation changed state and has
result data. Keep one aggregated wait mechanism, but define distinct object
semantics.

\subsection{Replace arbitrary cross-LP closures}

Boxed \texttt{FnOnce()} work transported through the IPI layer is hard to
account, difficult to inspect, and brittle under backpressure.
\textbf{Direction:} use closed kernel enums for mandatory kernel RPCs and
typed mailbox instances for subsystem-specific work.

\subsection{Do not equate shard wake with IPI delivery}

Remote wake mapping directly to \texttt{send\_ipi(target\_lp)} can generate
excessive interrupts. \textbf{Direction:} queue first, mark pending, send
a coalesced reschedule IPI only when the target cannot otherwise observe the
transition.

\subsection{Introduce memory objects before rich IPC payloads}

Xous's most important contribution --- move/lend semantics --- cannot be
implemented safely with transient pointers. \textbf{Direction:} prioritize
memory-object capability, mapping/unmapping, move between domains, temporary
lending, revocation on reply/cancel/death, and optional DMA integration.
Only then promote rich IPC payloads.

\subsection{Treat service discovery as userspace policy}

Do not add string-based global service lookup to the kernel. Build a userspace
name/policy service that receives bootstrap authority and delegates connection
capabilities.

\endinput
